\section*{Chapter \ref{ch:g2:water-on-earth} --- Water on Earth}

\begin{solution}{exo:g2:water-on-earth:1}
\omterm{def:g2:water-on-earth:salt-fresh}{Salt water}: the \omterm{def:g2:water-on-earth:ocean}{sea}, the \omterm{def:g2:water-on-earth:ocean}{ocean}. \omterm{def:g2:water-on-earth:salt-fresh}{Fresh water}: the \omterm{def:g1:landscapes-around-us:water}{river}, the rain, the \omterm{def:g1:landscapes-around-us:water}{lake}
in the \omterm{def:g1:landscapes-around-us:landform}{mountains}, the ice on the \omterm{def:g1:landscapes-around-us:landform}{mountain}.
\end{solution}

\begin{solution}{exo:g2:water-on-earth:2}
Because most of its surface is covered by the blue water of the \omterm{def:g2:water-on-earth:ocean}{oceans}:
seen from space, it looks blue.
\end{solution}

\begin{solution}{exo:g2:water-on-earth:3}
$100 - 71 = 29$ squares are land.
\end{solution}

\begin{solution}{exo:g2:water-on-earth:4}
An \omterm{def:g2:water-on-earth:ocean}{ocean} is bigger: a \omterm{def:g2:water-on-earth:ocean}{sea} is a smaller stretch of \omterm{def:g2:water-on-earth:salt-fresh}{salt water}, often at the
edge of an \omterm{def:g2:water-on-earth:ocean}{ocean}.
\end{solution}

\begin{solution}{exo:g2:water-on-earth:5}
No: the water of the \omterm{def:g2:water-on-earth:ocean}{ocean} is \omterm{def:g2:water-on-earth:salt-fresh}{salt water}, and drinking it makes the body
lose water. The sailor has to carry \omterm{def:g2:water-on-earth:salt-fresh}{fresh water}.
\end{solution}

\begin{solution}{exo:g2:water-on-earth:6}
For example: in \omterm{def:g1:landscapes-around-us:water}{rivers}, in \omterm{def:g1:landscapes-around-us:water}{lakes}, in the ice and snow of the \omterm{def:g1:landscapes-around-us:landform}{mountains} and
of Antarctica, under the ground, in the rain.
\end{solution}

\begin{solution}{exo:g2:water-on-earth:7}
\omterm{def:g2:water-on-earth:salt-fresh}{Salt water}: 96 buckets. \omterm{def:g2:water-on-earth:salt-fresh}{Fresh water}: $2 + 2 = 4$ buckets.
(Strictly, part of the water under the ground is salty too; the picture
counts it with the \omterm{def:g2:water-on-earth:salt-fresh}{fresh water}.)
\end{solution}

\begin{solution}{exo:g2:water-on-earth:8}
The \omterm{def:g2:water-on-earth:ocean}{oceans} (96 buckets), then the ice and the snow (2 buckets), then the
\omterm{def:g1:landscapes-around-us:water}{rivers} (much less than one bucket).
\end{solution}

\begin{solution}{exo:g2:water-on-earth:9}
It can soak into the ground; it can run over the ground into a stream or
a \omterm{def:g1:landscapes-around-us:water}{river}; it can evaporate back into the air when the sun comes out.
\end{solution}

\begin{solution}{exo:g2:water-on-earth:10}
Salt is left at the bottom of the pan. The water has evaporated: it has
gone into the air as water vapour.
\end{solution}

\begin{solution}{exo:g2:water-on-earth:11}
For example: the drop runs down the roof into the gutter and the drain,
then into a stream and a \omterm{def:g1:landscapes-around-us:water}{river}; the \omterm{def:g1:landscapes-around-us:water}{river} carries it to the \omterm{def:g2:water-on-earth:ocean}{sea}; one sunny
day it evaporates from the \omterm{def:g2:water-on-earth:ocean}{sea}, rises into the air and becomes part of a
cloud.
\end{solution}
